\chapter{Extended Technical Appendix}

This appendix provides additional implementation depth for the TCP/IP-inspired
UDP project. It is written to support evaluation, viva discussion, and future
continuation of the codebase.

\section{Appendix Purpose}
The main chapters cover the report view. This appendix keeps the practical
details, debug notes, and next-step ideas in one place.

\section{Complete Build and Run Workflow}
The project uses plain GCC and does not depend on external networking
frameworks. Each version follows the same basic run order.

\subsection{Compile}
For version $N$:
\begin{enumerate}
    \item \texttt{gcc -std=c99 -o receiver\_N receiver\_N.c}
    \item \texttt{gcc -std=c99 -o sender\_N sender\_N.c}
\end{enumerate}

\subsection{Execute}
\begin{enumerate}
    \item Open terminal 1 and run \texttt{./receiver\_N}
    \item Open terminal 2 and run \texttt{./sender\_N}
    \item Type the test message in the sender terminal.
\end{enumerate}

\subsection{Operational Notes}
\begin{enumerate}
    \item Start the receiver first so the bind path is ready.
    \item Use only one sender and receiver pair on port 9090 at a time.
    \item If a process hangs, stop the old binaries and run again.
\end{enumerate}

\section{Debug Checklist Used in This Project}
The following checklist was used repeatedly and is useful for future maintenance.

\begin{enumerate}
    \item Check the sender address and port.
    \item Check that the receiver binds cleanly and handles errors.
    \item Make sure packet fields are set before checksum work begins.
    \item Make sure checksum recomputation starts from a zero checksum field.
    \item Make sure sender base and next\_to\_send stay in order.
    \item Make sure ACK handling moves the window only once.
    \item Make sure duplicate packets are not delivered twice.
    \item Make sure the timeout is reasonable and retries cannot loop forever.
    \item Make sure the FIN exchange finishes and sockets close cleanly.
\end{enumerate}

\clearpage
\section{Detailed Behavior: Sender Core Loop}
Across advanced versions, sender logic follows a common event-driven pattern:

\begin{enumerate}
    \item Set up the socket, destination, and timeout values.
    \item Do the handshake when the level needs it.
    \item Split the message into packets or chunks.
    \item Send packets while keeping the in-flight count inside the window.
    \item Wait for ACKs and move the window forward.
    \item Watch for timeout or duplicate ACK events and retransmit when needed.
    \item Close the connection with FIN logic.
\end{enumerate}

\subsection{State Variables}
The most influential sender variables in high-level versions are:

\begin{enumerate}
    \item \texttt{base}: the oldest packet still waiting for ACK.
    \item \texttt{next\_to\_send}: the next packet ready to go out.
    \item \texttt{window\_size}: the current send limit.
    \item \texttt{rto\_ms}: the retransmission timeout in milliseconds.
    \item \texttt{cwnd} and \texttt{ssthresh}: the congestion-control values.
    \item duplicate ACK counters for each ACK number.
\end{enumerate}

\clearpage
\section{Detailed Behavior: Receiver Core Loop}
Advanced receiver implementations use predictable stages:

\begin{enumerate}
    \item Check the packet checksum and flags.
    \item Compare the sequence with \texttt{expected\_seq}.
    \item Deliver in-order data right away and move forward.
    \item Buffer out-of-order data that still fits in the window, then ACK it.
    \item Re-ACK duplicates so the sender can recover.
    \item Handle FIN and free resources cleanly.
\end{enumerate}

\subsection{Receiver Buffering in Selective-Repeat Style}
The selective-repeat style receiver uses parallel arrays for packets, lengths,
and presence. That keeps reordered data safe and lets delivery resume in the
right order once the missing packet arrives.

\section{Version-by-Version Pitfall Notes}
This section captures practical issues found during progressive development.

\subsection{v0--v2 Pitfalls}
\begin{enumerate}
    \item Bit order mistakes can rebuild the wrong character.
    \item Missing null terminator logic can leave the receiver waiting.
    \item \texttt{usleep} timing depends on the system.
\end{enumerate}

\subsection{v3--v7 Pitfalls}
\begin{enumerate}
    \item Parity checks need the same rules on both sides.
    \item ACK/NAK formats must be clear or the sender can stall.
    \item Retry loops need a limit so they do not run forever.
\end{enumerate}

\subsection{v8--v10 Pitfalls}
\begin{enumerate}
    \item Handshake flags must not clash.
    \item Recompute the checksum only after all payload bytes are set.
    \item Sliding-window pointer math can easily slip by one.
\end{enumerate}

\subsection{v11--v14 Pitfalls}
\begin{enumerate}
    \item Out-of-order flush must keep the sequence intact.
    \item Adaptive RTO can swing too much without steady RTT samples.
    \item Duplicate ACK counters must reset after each state change.
    \item Fast recovery can go wrong if ACKs are read too loosely.
\end{enumerate}

\section{Experiment Design Template}
A repeatable experiment protocol is essential for fair comparison. The template
used in this project is shown below.

\begin{table}[h]
\centering
\captionabove{Standard experiment template used for each protocol version.}
\renewcommand{\arraystretch}{1.18}
\begin{tabular}{|p{3.5cm}|p{9.4cm}|}
\hline
\textbf{Field} & \textbf{Description} \\
\hline
Version & sender\_N / receiver\_N pair \\
\hline
Message length & Short (20 chars), medium (60 chars), long (95 chars) \\
\hline
Fault setting & No fault / delay / simulated loss / duplicate packet scenario \\
\hline
Runs per case & Minimum 10 runs for stable average \\
\hline
Metrics & Completion time, retransmissions, ACK count, success/failure \\
\hline
Notes & Any abnormal behavior or timeout event logs \\
\hline
\end{tabular}
\label{tab:experiment-template}
\end{table}

\clearpage
\section{Sample Observation Narrative}
The following narrative style was used to explain the test runs.

\begin{displayquote}
In level 9, with a moderate induced delay, the sender kept timing out at the
window base. Retransmission from base brought progress back, but the overall
time to completion was much longer than in a run without delays. The receiver's
ACK pattern confirmed the overall behavior.
\end{displayquote}

\begin{displayquote}
In level 12, with a variable delay, adaptive RTO cut down on false timeout
events after the first RTT stabilized. The retransmission ratio went down when
compared to the fixed-timeout setup.
\end{displayquote}

\begin{displayquote}
In level 14, with a fake packet gap, the triple duplicate ACK condition caused
fast retransmit to happen before the timeout ran out. Recovery was clearly
faster than the path that was based on timeouts.
\end{displayquote}

\section{Illustrative State Machine View}
\vspace{2\baselineskip}
\begin{figure}[H]
\centering
\resizebox{\textwidth}{!}{%
\begin{tikzpicture}[
    stA/.style={draw, rounded corners=4pt, thick, minimum width=4.3cm, minimum height=1.15cm, align=center, fill=cyan!14, draw=cyan!55!black},
    stB/.style={draw, rounded corners=4pt, thick, minimum width=4.3cm, minimum height=1.15cm, align=center, fill=yellow!18, draw=orange!65!black},
    stC/.style={draw, rounded corners=4pt, thick, minimum width=4.3cm, minimum height=1.15cm, align=center, fill=green!14, draw=green!45!black},
    arr/.style={-{Latex[length=3.2mm]}, very thick, draw=black!82},
    lab/.style={fill=white, inner sep=1.2pt, font=\footnotesize}
]
\node[stA] (idle) at (0,0) {Idle};
\node[stB] (syn) at (5.8,0) {Handshake};
\node[stC] (send) at (11.6,0) {Send Window};
\node[stC] (wait) at (17.4,0) {Wait ACK};
\node[stA] (fin) at (23.2,0) {Teardown};
\node[stB] (retx) at (14.5,-2.9) {Fast Retransmit};

\draw[arr] (idle.east) -- (syn.west);
\draw[arr] (syn.east) -- (send.west);
\draw[arr] (send.east) -- (wait.west);
\draw[arr] (wait.east) -- node[above,lab]{all data acked} (fin.west);
\draw[arr] (wait.south) -- node[right,lab]{timeout / dup ACKs} (retx.north);
\draw[arr] (retx.west) to[bend right=18] node[below,lab]{retry path} (send.south);
\end{tikzpicture}%
}%

\caption{Simplified sender-side operational state flow in advanced versions.
}
\label{fig:appendix-state-machine}
\end{figure}

\section{Protocol Trace Formatting Recommendation}
\vspace{0.8\baselineskip}

\begin{figure}[H]
\centering
\resizebox{\textwidth}{!}{%
\begin{tikzpicture}[
    fld/.style={draw, rounded corners=4pt, thick, minimum width=4.6cm, minimum height=1.25cm, align=center, fill=teal!10, draw=teal!60!black, font=\small},
    eventfld/.style={draw, rounded corners=4pt, thick, minimum width=6.1cm, minimum height=1.25cm, align=center, fill=teal!10, draw=teal!60!black, font=\small},
    arr/.style={-{Latex[length=3.1mm]}, very thick, draw=black!82}
]
\node[fld] (f1) at (0,0) {[timestamp]};
\node[fld] (f2) at (5.5,0) {[role=sender|receiver]};
\node[eventfld] (f3) at (12.2,0) {[event=send|recv|ack|timeout|retransmit]};

\node[fld] (f4) at (0,-2.6) {[seq=x] [ack=y]};
\node[fld] (f5) at (5.5,-2.6) {[cwnd=z] [rto=ms]};
\node[fld] (f6) at (12.2,-2.6) {[note=text]};

\draw[arr] (f1.east) -- (f2.west);
\draw[arr] (f2.east) -- (f3.west);
\draw[arr] (f1.south) -- (f4.north);
\draw[arr] (f2.south) -- (f5.north);
\draw[arr] (f3.south) -- (f6.north);
\draw[arr] (f4.east) -- (f5.west);
\draw[arr] (f5.east) -- (f6.west);
\end{tikzpicture}%
}%

\caption{Recommended structured protocol-trace fields for automated analysis.
}
\label{fig:trace-format-diagram}
\end{figure}

\section{Suggested Quantitative Extensions}
Future project groups can convert this implementation into a mini benchmark
suite by adding automated scripts. Useful extensions include:

\begin{enumerate}
    \item throughput versus message size for each version,
    \item retransmission ratio versus injected packet loss,
    \item completion time under delay jitter,
    \item fairness with two parallel sender flows,
    \item adaptive RTO sensitivity to \(\alpha\) and \(\beta\).
\end{enumerate}

\section{Security and Reliability Perspective}
The project intentionally focuses on reliability and congestion behavior, not
security. But in real-world use, integrity and authenticity still matter.
Checksum catches random corruption, but it does not stop a malicious change.
A production extension would include:

\begin{enumerate}
    \item message authentication,
    \item optional encryption,
    \item replay protection,
    \item session identity checks.
\end{enumerate}

These sit outside the transport logic, but they matter in real systems.

\clearpage
\section{Code Organization Recommendations}
As code complexity grows, modularization becomes important. Recommended
refactoring pattern:

\begin{enumerate}
    \item \texttt{packet.h/.c} for packet structs and checksum helpers,
    \item \texttt{timer.h/.c} for RTT, RTO, and timestamp helpers,
    \item \texttt{window.h/.c} for send and receive window work,
    \item \texttt{congestion.h/.c} for cwnd and recovery steps,
    \item \texttt{main\_sender.c} and \texttt{main\_receiver.c} for the main flow.
\end{enumerate}

This separation reduces duplication across version files.

\section{Reproducibility Notes}
To ensure reproducibility of results in academic evaluation:

\begin{enumerate}
    \item note the compiler version and build flags,
    \item note the operating system and kernel details,
    \item close other network-heavy applications,
    \item run each experiment several times,
    \item report the average and spread (minimum/maximum or standard deviation).
\end{enumerate}

Without reproducibility controls, protocol comparisons may appear inconsistent.

\section{Expanded Testing Scenarios and Interpretation}
Some suggested scenarios are clean-channel baseline, delayed ACK, and packet loss.
injection, messages coming in out of order, duplicate ACK bursts, and messages that are almost too big
tests. These cases show correctness, recovery speed, timeout behavior, and
quality of congestion adaptation.

\section{Evaluation Rubric for Academic Assessment}
The following rubric can be used during project review.

\begin{table}[h]
\centering
\captionabove{Suggested assessment rubric for implementation quality.}
\renewcommand{\arraystretch}{1.18}
\begin{tabular}{|p{3.3cm}|>{\centering\arraybackslash}p{2.1cm}|p{7.3cm}|}
\hline
\textbf{Criterion} & \textbf{Weight} & \textbf{Evaluation Focus} \\
\hline
Correctness & 30\% & Data integrity checks, ordered delivery, clean teardown \\
\hline
Reliability & 20\% & Behavior under loss, timeout and retransmission policy quality \\
\hline
Efficiency & 15\% & Throughput trend improvement from low to high levels \\
\hline
Code Quality & 15\% & Clarity, modularity, safety checks, variable naming \\
\hline
Analysis & 10\% & Metric reporting and explanation quality \\
\hline
Documentation & 10\% & Diagram quality, chapter structure, reproducibility notes \\
\hline
\end{tabular}
\label{tab:assessment-rubric}
\end{table}

\clearpage
\section{Final Reflection}
From an educational standpoint, the strongest part of this project is the
clear link between theory and practice. Ideas from networking class turn into
behaviors you can see in terminal traces. The staged versions work like a
small course in protocol engineering, covering encoding, integrity checks,
reliability loops, pipelining, adaptation, and congestion response.

This project also shows that protocol design is about trade-offs. Better
reliability adds state. Higher throughput adds complexity. Faster recovery
needs careful event handling. Better adaptability needs steady measurement.
The main skill learned here is how to balance those trade-offs.

\section{Appendix Summary}
This appendix adds practical details, implementation tips, and ready-to-use
steps to the main report. Together with the main chapters, it tells the full
story from basic UDP communication to TCP-like reliability and congestion
control behavior in C. It is meant as a simple transport-control learning aid.
